% Post-checkpoint-14 research. Checked auxiliary statements; not in the CP14 PDF.
\begin{proposition}[Proved: contracting the zero-cost exceptional blocks]
\label{D:contract-zero-cost-blocks}
In the algebraic surface-germ setting of
Proposition~\ref{D:exceptional-degree-linear-core}, each connected
noncore block can be contracted to a rational double point.  The
contractions may be performed simultaneously, give a normal scheme
\(\overline Z\), and the morphism \(f:Z\to Y\) descends to
\(\overline Z\to Y\).  Away from the blocks the contraction is an
isomorphism.  The number of resulting rational double points is at
most \(2nc+n\), with the notation of that proposition.
\end{proposition}
\begin{proof}
Every component of a noncore block is a smooth rational curve of
square \(-2\); its intersection matrix is negative definite.
For every nonzero effective integral cycle \(T\) supported on the
block, adjunction gives
\[
 K_Z\cdot T=0,\qquad T^2\in2\mathbf Z_{<0},\qquad
                    p_a(T)=1+T^2/2\le0.                 \tag{D.RC.1}
\]
We apply Artin's contraction criterion
\citep[Theorem~2.3 and its full proof; published Russian
translation, pp.~9--10]{Artin1962D}.  Its hypotheses are precisely
negative definiteness and the inequalities in (D.RC.1); there is no
restriction on the characteristic.  We may first place the smooth
surface \(Z\) in a smooth projective compactification.  Indeed a
proper compactification can be resolved away from its smooth open
\(Z\), and a proper regular surface is projective
\citep[Lemma~54.16.11, tag~0C5P]{StacksDSurfaceResolution}.
Artin's theorem contracts the block on this projective surface to a
normal projective surface, changing nothing outside the block.
Restricting to the image of \(Z\) gives a scheme contraction of the
required open surface.  Distinct blocks are disjoint, so these
contractions can be carried out successively without changing the
remaining blocks or their intersection numbers.
For a cover of the punctured strict henselization of the algebraic
germ, first descend the finite cover, its resolution and the finite
list of indicated curves to an etale neighbourhood of finite type,
as in Proposition~\ref{D:surface-uniform-different}; perform this
construction there and then base change.  We do not use a projective
compactification of an arbitrary complete local nonalgebraic surface.

The resolution of each new singularity has precisely the contracted
smooth rational \((-2)\)-curves and has no exceptional curve of the
first kind.  The characteristic-free criterion for a rational
double point therefore applies
\citep[Proposition~0.4.8 and proof, author's April~2024 text,
p.~69]{C_Enriques2025}.  More explicitly, (D.RC.1) is Artin's
rationality criterion; for the fundamental cycle its arithmetic
genus is zero, hence its square is \(-2\), which gives multiplicity
two.  This is also the argument in the cited proof.

Each contracted block is connected and all its components map to
points of \(Y\).  Their images are therefore one point.  The graph
image of \(Z\to\overline Z\times Y\) is proper over \(\overline Z\)
and has one-point fibres.  It is thus finite and birational over
the normal surface \(\overline Z\), hence isomorphic to it.
This proves descent of \(f\).  The numerical bound is the block
bound of Proposition~\ref{D:exceptional-degree-linear-core}.
\end{proof}

\begin{lemma}[Proved: a uniform generator bound at rational double points]
\label{D:RDP-uniform-local-generators}
Let \(R\) be the complete local ring of a rational double point
over an algebraically closed field of positive characteristic.
Then
\[
          d\bigl(\pi_1(\operatorname{Spec}R-\{\mathfrak m\})\bigr)
                    \le6.                               \tag{D.RC.2}
\]
In particular its homomorphisms to any finite group \(H\) number
at most \(|H|^6\).  The bound is independent of Dynkin length,
characteristic and coindex, and does not assume a tame quotient
presentation of the singularity.
\end{lemma}
\begin{proof}
We use the complete table of local etale groups in
\citet[Convention~1.3(2), p.~81, and Table~7, p.~93]
{C_Matsumoto2022RDP}, recording Artin's calculations.
The unbounded-order families in that table are cyclic, dihedral,
and binary dihedral groups.  They have at most two generators;
the binary dihedral presentation is
\[
 \langle a,b\mid a^{2m}=1,\ b^2=a^m,\ bab^{-1}=a^{-1}\rangle.
\]
Every other listed group has order at most \(120\).  A finite
group of that order has at most six generators: in an irredundant
successive generating sequence, each proper subgroup enlargement
at least doubles the order, while \(2^7>120\).
Images of six fixed generators determine a homomorphism, proving
the last assertion.  The source specifically includes the
characteristic-two non-taut \(D\)-types; no group-scheme quotient
description is used.
\end{proof}

\begin{example}[Proved: a remaining branch invalidates a punctured-local bound]
\label{D:RDP-branch-gluing-obstruction}
Let \(k\) be algebraically closed of characteristic two and put
\[
 R=k[[u,v,w]]/(uv-w^3),\quad
 X=\operatorname{Spec}R,\quad x=\mathfrak m_R,\quad
 D=V(u,w).
\]
The full puncture \(X-\{x\}\) has fundamental group \(C_3\),
but
\[
               H^1_{\mathrm{et}}(X-D,\mathbf F_2)
                        \text{ is infinite}.             \tag{D.RC.3}
\]
Thus a uniform bound for the punctured group of a contracted
rational double point does not bound the covers of its local
complement when other removed curves pass through the point.
\end{example}
\begin{proof}
The ring \(R\) is the invariant ring of \(k[[s,t]]\) under
\(C_3\), acting by \((s,t)\mapsto(\zeta s,\zeta^{-1}t)\).
The invariant monomials are generated by
\(u=s^3,v=t^3,w=st\), with relation \(uv=w^3\).
The action is free away from the origin.  The puncture of
\(\operatorname{Spec}k[[s,t]]\) is simply connected by purity
\citep[Lemma~58.21.3(1), tag~0BMA]{StacksDLocalPurity}, since
every finite etale cover extends over the strictly henselian
regular local scheme.  The connected etale \(C_3\)-cover of
punctures therefore proves the first assertion.  Equivalently it
is the \(A_2\) row of the table used above.

We have \(X-D=D(u)\): modulo \(u\), the equation forces
\(w^3=0\).  For every \(a\in k\), consider the Artin--Schreier
torsor
\[
                   z^2-z=a/u.                           \tag{D.RC.4}
\]
It is finite etale on \(D(u)\).  At the height-one prime
\((u,w)\), the element \(v\) is a unit and \(w\) is a
uniformizer, so the normalized valuation of \(u\) is three.
For \(a\ne0\), the right side of (D.RC.4) has valuation \(-3\).
If \(b\) has negative valuation then
\(v(b^2-b)=2v(b)\) is even; if \(b\) has nonnegative valuation
then so does \(b^2-b\).  Thus \(a/u\) is not an
Artin--Schreier coboundary even in the fraction field.  Applying
this to \(a-a'\) shows that distinct constants give distinct
classes.  In fact (D.RC.4) defines an injection of additive
\(\mathbf F_2\)-spaces \(k\hookrightarrow H^1(D(u),\mathbf F_2)\).
This proves (D.RC.3).

The exact certificate
\nolinkurl{parts/C/code/check_RDP_branch_obstruction.py}
checks the finite monomial, group-action and valuation calculations;
its output is
\nolinkurl{parts/C/output/RDP_branch_obstruction.json}.
The infinite-family argument is the valuation proof just given.
\end{proof}

\begin{remark}[Exact obstruction to compression]
In the core construction the images of the core components usually
pass through the new rational double points.  The relevant cover is
then defined on the local complement of those branches.  It need
not extend to the full puncture, so (D.RC.2) cannot be inserted as a
bound for that local gluing problem.  The example is only a failed
inference for this method: it is not a counterexample to either
original finite-generation question.
\end{remark}
